\documentclass[a4paper,12pt]{article}
\usepackage[fontset=none]{ctex}
\usepackage{geometry}
\usepackage{graphicx}
\usepackage{amsmath}
\usepackage{enumitem}
\usepackage{setspace}
\usepackage{fontspec}

% 字体设置（与 docx 模板一致）
\setCJKmainfont{Songti SC}          % 宋体 - 正文默认
\setCJKsansfont{Heiti SC}            % 黑体 - 一级标题
\setCJKmonofont[AutoFakeBold]{STFangsong}          % 仿宋 - 注意事项
\setmainfont{Times New Roman}        % 英文默认

% 页面设置
\geometry{left=2.5cm, right=2.5cm, top=2.5cm, bottom=2.5cm}

% 去掉页码
\pagestyle{empty}

\begin{document}

% ========== 封面 ==========
\begin{center}
    \vspace*{2cm}
    {\fontsize{36pt}{43pt}\selectfont \textbf{物理实验预习报告}}
    \vspace{3cm}
\end{center}

\begin{center}
    \fontsize{16pt}{24pt}\selectfont
    \textbf{实验名称：弗兰克-赫兹实验\quad} \\[0.5em]
    \textbf{指导教师：\_\_\_\_\_\_\_\_\_\_\_\_\_\_\_\_\_\_\_\_\_\_\_\_\_}
\end{center}

\vspace{2cm}

\begin{center}
    \fontsize{14pt}{28pt}\selectfont
    \textbf{周次：\_\_\_\_\_\_\_\_} \\
    \textbf{姓名：\_\_\_\_\_\_\_\_} \\
    \textbf{学号：\_\_\_\_\_\_\_\_} \\[0.5em]
    \textbf{实验日期: \_\_\_\_年\_\_\_\_月\_\_\_\_日\quad 星期\_\_\_上/下午}
\end{center}

\vspace{1.5cm}

\begin{center}
    \fontsize{14pt}{20pt}\selectfont
    浙江大学物理实验教学中心
\end{center}

\newpage

% ========== 正文 ==========

{\fontsize{14pt}{20pt}\selectfont \textbf{1. 实验综述}}

\vspace{0.3cm}

{\fontsize{12pt}{18pt}\selectfont

弗兰克-赫兹实验是近代物理学中验证原子能级量子化的经典实验。1914年弗兰克与赫兹用慢电子轰击汞蒸气，测定了汞原子的第一激发电位，有力证明了原子内部存在分立的定态能级，为玻尔原子理论提供了直接实验证据，两人因此获得1925年诺贝尔物理学奖。

实验原理：在充氩弗兰克-赫兹管中，阴极$K$发射的热电子经加速电压$U_{G_2K}$加速后与氩原子碰撞。当电子能量小于氩原子第一激发能量时，仅发生弹性碰撞，板流$I_A$随$U_{G_2K}$增加而增大；当$U_{G_2K}$达到第一激发电势$U_0$时，电子与氩原子发生\textbf{非弹性碰撞}，将能量转移给氩原子使其跃迁至激发态，电子因失去能量无法克服拒斥电压$U_{G_2A}$到达板极，板流$I_A$下降。随$U_{G_2K}$继续增大，板流呈现周期性起伏，相邻峰（谷）值间的电势差即为氩原子第一激发电势，公认值约为$11.61$\,V。

实验装置为THQFH-2A弗兰克-赫兹实验仪，包括弗兰克-赫兹管（充氩四极管）、电源、微电流测量仪和示波器，通过自动扫描或手动逐点测量$I_A$--$U_{G_2K}$曲线，用逐差法求取第一激发电势。

}

\vspace{0.5cm}

{\fontsize{14pt}{20pt}\selectfont \textbf{2. 实验重点}}

\vspace{0.3cm}

{\fontsize{12pt}{18pt}\selectfont

理解原子能级量子化的物理图像及电子与原子碰撞激发的基本方法；掌握弗兰克-赫兹管各电极电压（灯丝电压$U_{F_1F_2}$、栅极电压$U_{G_1K}$、加速电压$U_{G_2K}$、拒斥电压$U_{G_2A}$）的作用及对$I_A$--$U_{G_2K}$曲线的影响；学会用逐差法从峰值电压数据中计算氩原子第一激发电势$U_0$。

}

\vspace{0.5cm}

{\fontsize{14pt}{20pt}\selectfont \textbf{3. 实验难点}}

\vspace{0.3cm}

{\fontsize{12pt}{18pt}\selectfont

合理选择灯丝电压$U_{F_1F_2}$、栅极电压$U_{G_1K}$及拒斥电压$U_{G_2A}$的参数组合，使$I_A$--$U_{G_2K}$曲线峰谷分明且不饱和；手动测量时在峰谷附近需加密测量点以准确定位极值；理解接触电势差导致第一个峰位偏移但不影响相邻峰间距的原因；正确运用逐差法处理数据并分析实验误差。

}

\vfill

% ========== 注意事项 ==========
\noindent\rule{\textwidth}{0.4pt}

{\CJKfontspec[AutoFakeBold]{STFangsong}\fontsize{12pt}{18pt}\selectfont
\textbf{注意事项：}
\begin{enumerate}[label=\arabic*., leftmargin=2em, nosep]
    \item 用PDF格式上传"预习报告"，文件名：学生姓名+学号+实验名称+周次。
    \item "预习报告"必须递交在"学在浙大"的本课程的对应实验项目的"作业"模块内。
    \item "预习报告"还须拷贝到"实验报告"中（便以教师批改）。
\end{enumerate}
}

\vspace{0.5cm}
\begin{center}
    {\CJKfontspec[AutoFakeBold]{STFangsong}\fontsize{12pt}{18pt}\selectfont \textbf{浙江大学物理实验教学中心制}}
\end{center}

\end{document}
